% Трёхфазный мостовой АИТ: источник U_d + большой дроссель L_d
% («источник тока»), мост VS1-VS6, нагрузки с C_к (звезда)
\documentclass[border=4pt]{standalone}
\usepackage{amsmath}
\usepackage{fontspec}
\setmainfont{Calibri}
\usepackage[european, RPvoltages]{circuitikz}
\definecolor{ink}{RGB}{31,42,51}
\definecolor{grp}{RGB}{46,110,140}
\ctikzset{bipoles/length=1.0cm}
\ctikzset{bipoles/thickness=0.9}
\ctikzset{diodes/fill=ink}
\ctikzset{inductor=american}

\begin{document}
\begin{tikzpicture}[line width=1.0pt, color=ink, font=\small]
  % --- источник U_d и дроссель L_d (источник тока) слева
  \draw (0,-0.6) to[V, l=$U_{\text{d}}$] (0,0.8) to[L, l=$L_{\text{d}}$] (0,4.6);
  \node[anchor=east] at (-0.15,0.85) {$+$};
  \node[anchor=east] at (-0.15,-0.55) {$-$};
  \draw[->, line width=0.9pt] (0.5,4.9) -- (1.3,4.9);
  \node[anchor=south] at (0.9,4.92) {$i_{\text{d}}$};
  % --- шины DC (объединяют катоды группы I и аноды группы II)
  \draw (0,4.6) -- (5.8,4.6);
  \draw (0,-0.6) -- (5.8,-0.6);
  % --- мост VS1-VS6: все остриём вверх, отводы фаз a, b, c
  \foreach \xc/\yph in {3.0/3.2, 4.4/2.2, 5.8/1.2}{
    \draw (\xc,\yph) -- (\xc,3.7) to[Ty] (\xc,4.3) -- (\xc,4.6) node[circ]{};
    \draw (\xc,-0.6) node[circ]{} -- (\xc,-0.1) to[Ty] (\xc,0.5) -- (\xc,\yph) node[circ]{};
  }
  \node[anchor=south, text=grp, font=\footnotesize] at (4.4,4.75)
       {катодная группа};
  \node[anchor=north, text=grp, font=\footnotesize] at (4.4,-1.25)
       {анодная группа};
  \node[anchor=south] at (2.35,4.68) {$+$};
  \node[anchor=north] at (2.35,-0.68) {$-$};
  % --- фазные провода к нагрузке
  \draw (3.0,3.2) -- (7.4,3.2) node[circ]{};
  \draw (4.4,2.2) -- (9.2,2.2) node[circ]{};
  \draw (5.8,1.2) -- (11.0,1.2) node[circ]{};
  \node[anchor=south west, font=\footnotesize] at (7.45,3.25) {a};
  \node[anchor=south west, font=\footnotesize] at (9.25,2.25) {b};
  \node[anchor=south west, font=\footnotesize] at (11.05,1.25) {c};
  % --- нагрузки: C_к параллельно Z_н (элементы ниже проводов фаз),
  %     звезда; шина нейтрали
  \draw (7.4,3.2) -- (7.4,0.7) to[C] (7.4,-0.3) -- (7.4,-1.0);
  \node[anchor=east] at (7.25,0.55) {$C_{\text{к}}$};
  \draw (7.4,3.2) -- (8.3,3.2) -- (8.3,0.7) to[generic] (8.3,-0.3) -- (8.3,-1.0);
  \draw (9.2,2.2) -- (9.2,0.7) to[C] (9.2,-0.3) -- (9.2,-1.0);
  \draw (9.2,2.2) -- (10.1,2.2) -- (10.1,0.7) to[generic] (10.1,-0.3) -- (10.1,-1.0);
  \draw (11.0,1.2) -- (11.0,0.7) to[C] (11.0,-0.3) -- (11.0,-1.0);
  \draw (11.0,1.2) -- (11.9,1.2) -- (11.9,0.7) to[generic, l=$Z_{\text{н}}$] (11.9,-0.3)
        -- (11.9,-1.0);
  \draw (7.4,-1.0) -- (11.9,-1.0);
  \node[circ] at (8.3,-1.0) {};
  \node[circ] at (9.2,-1.0) {};
  \node[circ] at (10.1,-1.0) {};
  \node[circ] at (11.0,-1.0) {};
  % --- подписи
  \node[anchor=west] at (3.35,4.0)  {VS1};
  \node[anchor=west] at (4.75,4.0)  {VS2};
  \node[anchor=west] at (6.15,4.0)  {VS3};
  \node[anchor=west] at (3.35,0.2)  {VS4};
  \node[anchor=west] at (4.75,0.2)  {VS5};
  \node[anchor=west] at (6.15,0.2)  {VS6};
  \node[anchor=south, text=grp, font=\footnotesize] at (9.6,-1.5)
       {$C_{\text{к}}$, $Z_{\text{н}}$ — в каждой фазе (звезда)};
  \node[anchor=south, text=grp, font=\footnotesize] at (2.2,5.4)
       {$L_{\text{d}}$ — источник тока в звене постоянного тока};
\end{tikzpicture}
\end{document}
